% Ten-minute talk, eight slides. Build: ./build.sh  (or pdflatex slides.tex twice)
% Every number on the slides is copied from manuscript/numbers.tex,
% short-summary/comparison-numbers.tex or prompts/uso-de-tiempo-y-tokens.md.
\documentclass[aspectratio=169,11pt]{beamer}
\usetheme[progressbar=frametitle,numbering=fraction]{metropolis}
\usepackage[T1]{fontenc}
\usepackage[utf8]{inputenc}
\usepackage{lmodern}
\usepackage{amsmath}
\usepackage{graphicx}
\usepackage{tikz}
\usetikzlibrary{positioning,arrows.meta}
\graphicspath{{figures/}}

\definecolor{ink}{HTML}{1F2937}
\definecolor{accent}{HTML}{1F5FA8}
\definecolor{warm}{HTML}{D9622B}
\definecolor{muted}{HTML}{6B7280}
\definecolor{paper}{HTML}{FFFFFF}
\setbeamercolor{normal text}{fg=ink,bg=paper}
\setbeamercolor{frametitle}{fg=white,bg=accent}
\setbeamercolor{progress bar}{fg=warm,bg=accent!30}
\setbeamercolor{alerted text}{fg=warm}
\setbeamerfont{frametitle}{size=\large}

\newcommand{\ext}{\mathrm{ext}}
\newcommand{\stat}[2]{{\large\bfseries\color{accent}#1\par}\vspace{1pt}{\scriptsize\color{muted}#2\par}}
\newcommand{\src}[1]{{\scriptsize\color{muted}#1}}
\newcommand{\quoted}[2]{%
  \begin{beamercolorbox}[wd=\linewidth,sep=5pt,leftskip=4pt]{block body}
  \footnotesize\itshape ``#1''\par\smallskip\upshape\scriptsize\color{muted}#2
  \end{beamercolorbox}\vspace{3pt}}
\setbeamercolor{block body}{bg=accent!7}

\title{Extremality shift of rotating black holes\\at finite Gauss--Bonnet coupling}
\subtitle{A research project carried out with AI agents}

\begin{document}

% ---------------------------------------------------------------- 1
\begin{frame}[plain]
\vspace*{0.4cm}
{\LARGE\bfseries\color{accent}Extremality shift of rotating black holes\\[2pt]
at finite Gauss--Bonnet coupling\par}
\vspace{0.25cm}
{\large A research project carried out with AI agents\par}
\vspace{0.55cm}
{\normalsize Cecilia Bejarano$^{1}$, David Blanco$^{2}$, Andrés Goya$^{1}$, Guillem Pérez-Nadal$^{2}$\par}
\vspace{0.12cm}
{\scriptsize\color{muted}$^1$Instituto de Astronomía y Física del Espacio (IAFE, CONICET--UBA)\quad
$^2$Departamento de Física, FCEN--UBA and IFIBA (CONICET--UBA)\par}
\vspace{0.3cm}
{\small With AI assistance: \textbf{Claude Code} (Anthropic; Claude Opus~5, Sonnet~5)
and \textbf{Codex} (OpenAI; GPT-5)\par}
\vfill
{\footnotesize\color{muted}
\emph{Gravitational Waves and AI-Assisted Research} --- M.~Zaldarriaga \& N.~Mirón-Granese\\
Departamento de Física, FCEN, Universidad de Buenos Aires \quad\textperiodcentered{}\quad 1--2 October 2026
\quad\textperiodcentered{}\quad \texttt{github.com/ddblanco/BHs}\par}
\end{frame}

% ---------------------------------------------------------------- 2
\begin{frame}{The question: how far does Gauss--Bonnet move the extremal bound?}
\begin{columns}[T,onlytextwidth]
\column{.52\textwidth}
\centering\includegraphics[width=.8\linewidth]{fig-question.pdf}

\vspace{2pt}
\begin{beamercolorbox}[wd=\linewidth,sep=3pt,leftskip=2pt]{block body}
\raggedright\scriptsize
\textbf{Warm-up: Kerr, 4D} {\color{muted}($G=c=1$, $a=J/M$)}\\[1pt]
Extremality: $J\le M^2$. At $J=M^2$ the horizons merge and
$T=\sqrt{M^2-a^2}/(4\pi M r_+)=0$.\\[1pt]
First law: $dM=T\,dS+\Omega_H\,dJ$, \ with $T=\kappa/2\pi$, $S=A/4$.
\end{beamercolorbox}
\column{.45\textwidth}
\small
\textbf{Einstein, 5D, two equal spins $J$:}
\[M\ \ge\ \tfrac32\,\pi^{1/3}J^{2/3}\]
saturated by an \alert{extremal} ($T=0$) black hole.

\medskip
\textbf{Add Gauss--Bonnet:}
\[I=\tfrac1{16\pi G}\!\int\!\sqrt{-g}\,\bigl(R+\alpha\,\mathcal L_{\rm GB}\bigr)\]
\vspace{-7pt}
\par\src{schematically \ $R\sim g\,\partial^2g+\partial g\,\partial g$,\quad $\mathcal L_{\rm GB}\sim R^2$}\par

\medskip
\textbf{Known:} slope $\pi$ at first order in $\alpha$\\
\src{Ma, Li, Lü, arXiv:2009.00015}

\medskip
\textbf{Asked:} the slope $\partial M_\ext/\partial\alpha|_J$\\ at \alert{finite} $\alpha$.
\end{columns}
\end{frame}

% ---------------------------------------------------------------- 3
\begin{frame}[t]{How it started: the first prompts}
\vspace{-4pt}
\begin{beamercolorbox}[wd=\linewidth,sep=4pt,leftskip=4pt]{block body}
\scriptsize\itshape ``El proyecto consiste en utilizar agentes de IA para implementar códigos
numéricos que permitan obtener soluciones de agujeros negros rotantes en
dimensión D>4 [\ldots] reproducir soluciones numéricas ya encontradas en la
literatura. [\ldots] el paso siguiente sería emprender la búsqueda de nuevas
soluciones.''\par\smallskip\upshape\tiny\color{muted}Prompt 1 --- 4 Sept, 19:16, to Codex
\quad\textperiodcentered{}\quad then ``Fijate que el proyecto realmente siga las indicaciones propuestas por el
docente del curso'' (3), ``ejecutar con subagentes'' (5), ``continuemos con el hito 3'', ``dale, avanzá''
\quad\textperiodcentered{}\quad \textbf{79} human prompts, 4--28 Sept \quad\textperiodcentered{}\quad Claude Code builds $\leftrightarrow$ Codex referees
\end{beamercolorbox}
\par\vspace{-16pt}\noindent
\begin{minipage}[t]{.44\textwidth}\vspace{0pt}\centering
\includegraphics[width=.9\linewidth]{fig-roadmap.png}\par\vspace{2pt}
{\tiny\color{muted}first roadmap, 5 Sept. \textcolor{warm}{Hito 6} became $\Psi$;
Hito 8 ($T\to0$, paper, referees) was not in the plan.\par}
\end{minipage}\hfill
\begin{minipage}[t]{.52\textwidth}\vspace{0pt}\centering
\includegraphics[width=\linewidth]{fig-smarr.png}\par\vspace{2pt}
{\tiny\color{muted}arXiv:1010.0860, \S4.1.3 and footnote 14.\par}
\vspace{4pt}
{\scriptsize Hito 6, $\alpha$ as a thermodynamic variable:\
$2M=3TS+6\Omega_HJ+2\alpha\Psi$\par}
\end{minipage}
\end{frame}

% ---------------------------------------------------------------- 4
\begin{frame}{Step 1: reproduce the published family (arXiv:1010.0860, fig.~1)}
\centering
\includegraphics[width=.7\textwidth]{fig-paper-comparison.pdf}

\vspace{2pt}
\begin{columns}[T,onlytextwidth]
\column{.32\textwidth}\centering
\stat{$2\times10^{-4}$}{agreement at the horizon\\(a third of their plotting resolution)}
\column{.36\textwidth}\centering
\stat{$\sim10^{-2}$}{interior gap --- but their curve is\\ as far from the \emph{exact} solution}
\column{.32\textwidth}\centering
\stat{$1.1021$ vs $1.104$}{ergosurface radius:\\ open, reported as such}
\end{columns}
\vspace{2pt}
\src{Curves read as vectors from the arXiv PostScript source, not digitised from the image.}
\end{frame}

% ---------------------------------------------------------------- 5
\begin{frame}{Step 2: measure the slope on each solution, not from differences}
\centering
\small
$\displaystyle dM=T\,dS+2\Omega_H\,dJ+\Psi\,d\alpha
\quad\Longrightarrow\quad
\left.\frac{\partial M_\ext}{\partial\alpha}\right|_J=\lim_{T\to0}\Psi$
\qquad
$\displaystyle \mathsf A\,\frac{du}{d\alpha}=-\frac{\partial R}{\partial\alpha}$
\src{\;(one linear solve)}

\vspace{2pt}
\begin{columns}[T,onlytextwidth]
\column{.5\textwidth}\centering
\includegraphics[width=.86\linewidth]{fig-profiles.pdf}\\
\footnotesize\textbf{356} nonlinear solutions, \textbf{13} couplings $0\le\alpha\le\tfrac12$
\column{.5\textwidth}\centering
\includegraphics[width=.86\linewidth]{fig-potential.pdf}\\
\footnotesize $\Psi$ against temperature; the slope is the value at $T=0$
\end{columns}
\end{frame}

% ---------------------------------------------------------------- 6
\begin{frame}{Result: the bound tightens, but not linearly}
\begin{columns}[T,onlytextwidth]
\column{.45\textwidth}\centering
\includegraphics[width=\linewidth]{fig-shift.pdf}\\
\footnotesize slope $\partial M_\ext/\partial\alpha|_J$
\column{.31\textwidth}\centering
\includegraphics[width=\linewidth,trim=0 0 208.1bp 0,clip]{fig-massext.pdf}\\
\footnotesize extremal mass
\column{.22\textwidth}\raggedright
\vspace{2pt}
\stat{$>0$}{everywhere}\vspace{6pt}
\stat{$\pi\to1.07$}{then up to $1.44$}\vspace{6pt}
\stat{$+31\%$}{$M_\ext$ growth}\vspace{6pt}
\stat{$\times1.32$}{first-order overshoot}
\end{columns}
\vspace{4pt}
\src{\centering $T=0$ values are extrapolations from the coldest states reached ($TJ^{1/3}\approx3\times10^{-6}$).\par}
\end{frame}

% ---------------------------------------------------------------- 7
\begin{frame}{How far can it be trusted?}
\begin{columns}[T,onlytextwidth]
\column{.6\textwidth}\centering
\includegraphics[width=\linewidth]{fig-checks.pdf}
\column{.37\textwidth}\small
\textbf{Independent checks}\\ none uses the linear-response solve.

\medskip
\textbf{Agents refereeing agents}\\ a requested check \emph{failed}; an exact bound caught a real mass error.

\medskip
\textbf{Still open}\\
extrapolation to $T=0$;\\ equal spins; one branch only.
\end{columns}
\end{frame}

% ---------------------------------------------------------------- 8
\begin{frame}{What it cost}
\centering
\begin{columns}[T,onlytextwidth]
\column{.25\textwidth}\centering\stat{36.8 h}{active agent time, 13 days}
\column{.25\textwidth}\centering\stat{1.42 billion}{tokens processed\\(99\% re-read context)}
\column{.25\textwidth}\centering\stat{5.9 million}{tokens written ($\approx$24 MB)}
\column{.25\textwidth}\centering\stat{79}{human prompts}
\end{columns}
\vspace{4pt}
\includegraphics[width=.8\textwidth]{fig-cost.pdf}\\
\small The first \alert{new} object cost $5\times$ the tokens of reproducing everything known.\\[2pt]
\src{Not counted: human time, background numerical runs.}
\end{frame}

\end{document}
